\textbf{Phần 1. Trắc nghiệm lựa chọn}
\setcounterpageref{ex}{0}

\begin{ex} %[3H2V21]
Một khối lập phương có bao nhiêu đỉnh, bao nhiêu cạnh và bao nhiêu mặt?
\choice
{$6$ đỉnh, $12$ cạnh, $8$ mặt}
{$8$ đỉnh, $6$ cạnh, $12$ mặt}
{\True $8$ đỉnh, $12$ cạnh, $6$ mặt}
{$12$ đỉnh, $8$ cạnh, $6$ mặt}
\loigiai{
Khối lập phương có 8 đỉnh, 12 cạnh bằng nhau và 6 mặt đều là hình vuông.
\textbf{Đáp án của câu hỏi là C}
}
\end{ex}

\begin{ex} %[3H2V21]
Các mặt của khối lập phương đều là hình gì?
\choice
{\True Hình vuông}
{Hình chữ nhật}
{Hình tam giác}
{Hình tròn}
\loigiai{
Khối lập phương có 6 mặt đều là các hình vuông bằng nhau.
\textbf{Đáp án của câu hỏi là A}
}
\end{ex}

\begin{ex} %[3H2V21]
Một con xúc xắc chuẩn hình lập phương có tổng số chấm ở hai mặt đối diện luôn luôn bằng
\choice
{$6$}
{$8$}
{$9$}
{\True $7$}
\loigiai{
Quy ước chuẩn của xúc xắc: 1 đối diện 6, 2 đối diện 5, 3 đối diện 4; tổng số chấm luôn bằng 7.
\textbf{Đáp án của câu hỏi là D}
}
\end{ex}

\begin{ex} %[3H2V21]
Khối hộp chữ nhật có tất cả bao nhiêu mặt?
\choice
{$4$ mặt}
{\True $6$ mặt}
{$8$ mặt}
{$12$ mặt}
\loigiai{
Khối hộp chữ nhật có 6 mặt (2 mặt đáy và 4 mặt bên).
\textbf{Đáp án của câu hỏi là B}
}
\end{ex}

\begin{ex} %[3H2V21]
Dùng các khối lập phương nhỏ cạnh $1\text{ cm}$ để xếp thành một khối lập phương lớn hơn có cạnh $2\text{ cm}$. Số khối lập phương nhỏ cần dùng là
\choice
{$4$ khối}
{$6$ khối}
{\True $8$ khối}
{$16$ khối}
\loigiai{
Mỗi lớp cần $2 \times 2 = 4$ khối, 2 lớp cần: $4 \times 2 = 8$ khối lập phương nhỏ.
\textbf{Đáp án của câu hỏi là C}
}
\end{ex}

\begin{ex} %[3H2V21]
Người ta dùng dây ruy-băng để buộc một hộp quà hình lập phương cạnh $10\text{ cm}$ theo hình chữ thập quanh hộp (nút thắt cần thêm $15\text{ cm}$ ruy-băng). Độ dài sợi dây ruy-băng cần dùng là
\choice
{\True $55\text{ cm}$}
{$40\text{ cm}$}
{$50\text{ cm}$}
{$60\text{ cm}$}
\loigiai{
Buộc chữ thập quanh 4 mặt bên và 2 đáy hết: $10 \times 4 = 40\text{ cm}$. Thêm nút thắt $15\text{ cm}$, tổng chiều dài: $40 + 15 = 55\text{ cm}$.
\textbf{Đáp án của câu hỏi là A}
}
\end{ex}

\begin{ex} %[3H2V21]
Một khối gỗ hình hộp chữ nhật có chiều dài $5\text{ cm}$, chiều rộng $3\text{ cm}$, chiều cao $2\text{ cm}$. Người ta cắt khối gỗ đó thành các khối lập phương nhỏ cạnh $1\text{ cm}$. Số khối lập phương nhỏ cắt được là
\choice
{$10$ khối}
{$20$ khối}
{$25$ khối}
{\True $30$ khối}
\loigiai{
Số khối nhỏ cắt được là: $5 \times 3 \times 2 = 30$ khối lập phương nhỏ.
\textbf{Đáp án của câu hỏi là D}
}
\end{ex}

\begin{ex} %[3H2V21]
Khung của một chiếc lồng bàn hình khối lập phương được làm bằng các nan tre uốn theo các cạnh của nó. Biết cạnh khối lập phương là $15\text{ cm}$. Tổng chiều dài các nan tre cần dùng là
\choice
{$120\text{ cm}$}
{\True $180\text{ cm}$}
{$150\text{ cm}$}
{$90\text{ cm}$}
\loigiai{
Khối lập phương có 12 cạnh bằng nhau. Tổng chiều dài các cạnh: $12 \times 15 = 180\text{ (cm)}$.
\textbf{Đáp án của câu hỏi là B}
}
\end{ex}

\begin{ex} %[3H2V21]
Sơn tất cả các mặt ngoài của một khối lập phương được ghép bởi 8 khối lập phương nhỏ (cạnh $2\text{ cm}$). Hỏi có bao nhiêu khối lập phương nhỏ được sơn đúng 3 mặt?
\choice
{$4$ khối}
{$6$ khối}
{\True $8$ khối}
{$0$ khối}
\loigiai{
Cả 8 khối nhỏ đều nằm ở 8 đỉnh của khối lập phương lớn nên mỗi khối nhỏ đều được sơn đúng 3 mặt ngoài.
\textbf{Đáp án của câu hỏi là C}
}
\end{ex}

\begin{ex} %[3H2V21]
Khẳng định nào sau đây là ĐÚNG khi nói về khối lập phương và khối hộp chữ nhật?
\choice
{\True Khối lập phương là một dạng đặc biệt của khối hộp chữ nhật có chiều dài, chiều rộng và chiều cao bằng nhau}
{Khối hộp chữ nhật có nhiều đỉnh hơn khối lập phương}
{Khối lập phương có 8 mặt}
{Khối hộp chữ nhật có 16 cạnh}
\loigiai{
Khối lập phương là trường hợp đặc biệt của khối hộp chữ nhật khi tất cả các kích thước bằng nhau.
\textbf{Đáp án của câu hỏi là A}
}
\end{ex}

\begin{ex} %[3H2V21]
Một khối lập phương lớn có cạnh bằng $3\text{ cm}$ được tạo bởi 27 khối lập phương nhỏ cạnh $1\text{ cm}$. Người ta lấy đi 1 khối nhỏ ở chính giữa mặt trên. Khối còn lại có tất cả bao nhiêu mặt vuông nhỏ cạnh $1\text{ cm}$ ở bề mặt ngoài?
\choice
{$54$ mặt}
{$56$ mặt}
{$50$ mặt}
{\True $58$ mặt}
\loigiai{
Ban đầu khối lớn có $6 \times 9 = 54$ mặt nhỏ. Khi bỏ 1 khối nhỏ ở giữa mặt trên, mất đi 1 mặt ngoài nhưng lộ ra 5 mặt bên trong hốc: $54 - 1 + 5 = 58$ mặt vuông nhỏ.
\textbf{Đáp án của câu hỏi là D}
}
\end{ex}

\begin{ex} %[3H2V21]
Một hình hộp chữ nhật có kích thước đáy là $4\text{ cm}$ và $6\text{ cm}$, chiều cao là $5\text{ cm}$. Diện tích của 2 mặt đáy cộng lại là
\choice
{$24\text{ cm}^2$}
{\True $48\text{ cm}^2$}
{$50\text{ cm}^2$}
{$60\text{ cm}^2$}
\loigiai{
Diện tích một mặt đáy là: $4 \times 6 = 24\text{ cm}^2$. Tổng diện tích 2 đáy là: $24 \times 2 = 48\text{ cm}^2$.
\textbf{Đáp án của câu hỏi là B}
}
\end{ex}

\textbf{Phần 2. Trắc nghiệm đúng sai}
\setcounterpageref{ex}{0}

\begin{ex} %[3H2V21]
Xét tính đúng sai của các phát biểu về khối lập phương và khối hộp chữ nhật:
\choiceTF
{\True Khối lập phương và khối hộp chữ nhật đều có đúng 8 đỉnh và 12 cạnh.}
{Tất cả các mặt của khối hộp chữ nhật luôn luôn phải là hình vuông.}
{\True Một khối lập phương có 6 mặt đều là các hình vuông bằng nhau.}
{Khối lập phương có 6 đỉnh và 8 mặt.}
\loigiai{
\begin{itemize}
    \item Ý a) Đúng: Số đỉnh và cạnh của hai khối là như nhau.
    \item Ý b) Sai: Sai vì các mặt thường là hình chữ nhật.
    \item Ý c) Đúng: Đặc tính định nghĩa khối lập phương.
    \item Ý d) Sai: Sai vì có 8 đỉnh và 6 mặt.
\end{itemize}
\begin{itemize}
    \item a) \textbf{Đúng}.
    \item b) \textbf{Sai}.
    \item c) \textbf{Đúng}.
    \item d) \textbf{Sai}.
\end{itemize}
}
\end{ex}

\begin{ex} %[3H2V21]
Một khối gỗ hình lập phương có cạnh dài $10\text{ cm}$:
\choiceTF
{\True Tổng độ dài của tất cả 12 cạnh của khối gỗ là $120\text{ cm}$.}
{\True Mỗi mặt của khối gỗ là một hình vuông có chu vi $40\text{ cm}$.}
{Nếu cưa khối gỗ thành hai nửa bằng một nhát cưa phẳng thì hai khối thu được là khối lập phương.}
{\True Hai mặt đối diện bất kì của khối gỗ luôn luôn bằng nhau.}
\loigiai{
\begin{itemize}
    \item Ý a) Đúng: $12 \times 10 = 120\text{ cm}$.
    \item Ý b) Đúng: $10 \times 4 = 40\text{ cm}$.
    \item Ý c) Sai: Sai vì hai khối thu được là hình hộp chữ nhật có cạnh $10, 10, 5\text{ cm}$.
    \item Ý d) Đúng: Các mặt của khối lập phương đều bằng nhau.
\end{itemize}
\begin{itemize}
    \item a) \textbf{Đúng}.
    \item b) \textbf{Đúng}.
    \item c) \textbf{Sai}.
    \item d) \textbf{Đúng}.
\end{itemize}
}
\end{ex}

\textbf{Phần 3. Trả lời ngắn}
\setcounterpageref{ex}{0}

\begin{ex} %[3H2V21]
Một khối lập phương có cạnh bằng $8\text{ cm}$. Tính tổng độ dài tất cả các cạnh của khối lập phương đó theo xăng-ti-mét.
\shortans{96}
\loigiai{
\textbf{Đáp số:} $96$\\
Khối lập phương có 12 cạnh: $12 \times 8 = 96\text{ (cm)}$.
}
\end{ex}

\begin{ex} %[3H2V21]
Để làm một khung đèn lồng hình hộp chữ nhật có kích thước dài $20\text{ cm}$, rộng $15\text{ cm}$ và cao $10\text{ cm}$, người ta cần dùng bao nhiêu xăng-ti-mét dây thép (bỏ qua mối nối)?
\shortans{180}
\loigiai{
\textbf{Đáp số:} $180$\\
Hình hộp chữ nhật có 4 chiều dài, 4 chiều rộng, 4 chiều cao. Tổng độ dài: $4 \times (20 + 15 + 10) = 4 \times 45 = 180\text{ (cm)}$.
}
\end{ex}

\begin{ex} %[3H2V21]
Cần bao nhiêu khối lập phương nhỏ cạnh $1\text{ cm}$ để ghép thành một khối lập phương lớn có cạnh $3\text{ cm}$?
\shortans{27}
\loigiai{
\textbf{Đáp số:} $27$\\
Số khối nhỏ cần dùng: $3 \times 3 \times 3 = 27$ khối.
}
\end{ex}

\begin{ex} %[3H2V21]
Một người xếp các viên gạch hình hộp chữ nhật thành một khối gồm 4 lớp, mỗi lớp có 6 viên gạch. Hỏi khối đó có tất cả bao nhiêu viên gạch?
\shortans{24}
\loigiai{
\textbf{Đáp số:} $24$\\
Tổng số viên gạch: $6 \times 4 = 24$ viên.
}
\end{ex}

\begin{ex} %[3H2V21]
Một khối lập phương có chu vi một mặt đáy là $28\text{ cm}$. Hỏi cạnh của khối lập phương đó dài bao nhiêu xăng-ti-mét?
\shortans{7}
\loigiai{
\textbf{Đáp số:} $7$\\
Mặt đáy là hình vuông nên cạnh là: $28 : 4 = 7\text{ (cm)}$.
}
\end{ex}

\begin{ex} %[3H2V21]
Một xúc xắc chuẩn hình lập phương có các mặt ghi từ 1 đến 6 chấm. Tính tổng số chấm trên tất cả 6 mặt của con xúc xắc đó.
\shortans{21}
\loigiai{
\textbf{Đáp số:} $21$\\
Tổng số chấm: $1 + 2 + 3 + 4 + 5 + 6 = 21$.
}
\end{ex}
